\section{A ground-truth-free protocol}
\label{sec:protocol}

\paragraph{The rule.}
The protocol takes the measurement $y$, a solver $S$ with a consistency threshold $\tau_S$, an orientation scorer $u$, a target $M$ and a budget $N_{\max}$ (\cref{fig:teaser}a).
(1)~Draw runs $\hat x_1,\hat x_2,\dots$ of $S$ with fresh seeds and stop as soon as $M$ of them have $\rho\le\tau_S$, or after $N_{\max}$ runs.
(2)~Keep $\hat x=\argmin_j\rho(\hat x_j)$.
(3)~Output $R_{111}\hat x$ if $\delta=u(R_{111}\hat x)-u(\hat x)>0$ and $\hat x$ otherwise.
(4)~Raise $\textsc{flag}_C$ if no drawn run had $\rho\le\tau_S$, and $\textsc{flag}_O$ if $|\delta|$ falls in the lowest $30\%$ of the margins of the evaluated outputs.
We use $M=2$ throughout.
Each step uses the measurement, the solver and a zero-shot scorer, never the ground truth.
The design choices follow from what the residual can and cannot tell.

\paragraph{The residual compares runs; it does not certify them.}
Within an image, a \good run has a smaller residual than an \incons run in $570$ of $571$ pairs, and than a \cons run in $200$ of $205$\claim{C-A7}.
The comparison is sharp because two runs of the same image are scored against the same noise, which cancels: a \cons run of SITCOM fits worse than a \good run of the same image by a median of only $0.7\%$ of $\sigma$, yet the order is systematic\claim{C-A8}.
Absolute residual levels, in contrast, are set by the solver.
With signed measurements the ground truth has $\rho(\xs)=\sigma$ to within $0.3\%$, but the \good runs of DAPS and SITCOM, whose data terms are overweighted (\cref{sec:setup}), have median residuals of $0.946\sigma$ and $0.884\sigma$, against $0.997\sigma$ for FAST-DIPS\claim{C-A6}.
A threshold at the noise level would accept most of their failed runs, and one fitted to their \good runs would reject the ground truth.
The protocol therefore uses $\tau_S$ only to decide when enough candidates exist and when none does, and $\textsc{flag}_C$ means that no run reached the solver's own consistency level, not that the output is wrong.

\paragraph{Calibrating $\tau_S$ once per solver.}
$\tau$ behaves as a constant of the solver and noise level: it is $0.0501$ for FAST-DIPS on FFHQ and $0.0502$ on ImageNet, and differs by less than $0.3\%$ between our two FFHQ cohorts for DAPS ($0.0475$, $0.0476$) and SITCOM ($0.0449$, $0.0450$)\claim{C-G3}.
We therefore calibrate it once on images with ground truth and use it unchanged elsewhere; on ImageNet the protocol uses the FFHQ value.
Where no calibration set exists, a label-free threshold, the median plus three scaled MADs of the per-image minimum residual over a pilot of four runs, lands within $2\%$ of $\tau$ for every solver and cohort\claim{C-G3}.

\paragraph{Orientation.}
By \cref{sec:setup}, the choice within an orbit must come from the prior, and any test that orders orbit members like the prior is the Bayes choice.
On FFHQ, a zero-shot CLIP test with fixed upright and upside-down face prompts fixes all $58$ global orbit errors among $1{,}960$ runs and flips none of the $1{,}171$ \good runs\claim{C-A2}.
On ImageNet the generic prompts ``a photo.'' and ``an upside-down photo.'' repair $57$ of $67$ global orbit errors of FAST-DIPS and flip $5$ of $100$ \good runs\claim{C-A11}, and the margin of the test is informative: among the development outputs, those outside the least confident $30\%$ are \good in $65.3\%$ of cases, against $50.7\%$ for all outputs without $\textsc{flag}_C$\claim{C-G4}.

\paragraph{What stopping buys, and when.}
Fixed best-of-$K$ spends $K$ runs on every image; the stopping rule spends few runs on an image whose runs are usually consistent and $N_{\max}$ on one whose runs never are, which it then flags.
\begin{proposition}[stopping at two consistent runs]
\label{prop:stop}
Fix an image whose runs pass $\rho\le\tau_S$ independently with probability $s$, let $r=1-s$, and let $H_j^+$ ($H_N^-$) be the probability that residual selection and the orientation test return a \good image from $j$ passing ($N$ non-passing) runs.
With $M=2$ and budget $N_{\max}=N$,
\begin{align*}
\E T&=\bigl(2-Nr^{N-1}+(N-2)r^N\bigr)/s,\\
\Prob(\textsc{flag}_C)&=r^N,\\
\Prob(\mathrm{success})&=(1-r^N-Nsr^{N-1})\,H_2^+\\
&\quad+Nsr^{N-1}H_1^+ + r^NH_N^-.
\end{align*}
\end{proposition}
The proof rests on one fact: the first two passing runs are independent draws from the passing law, whenever they arrive (supplement)\claim{C-T12}.
The consistency rate $s$ thus decides how many runs are spent, and $H^+$, how often a consistent run is correct, decides what they buy.
Heterogeneity of $s$ across images does not by itself make stopping better than a fixed $K$ at equal mean compute; the supplement gives a counterexample.
Stopping gains when consistency predicts correctness: when an image's failed runs are \incons rather than \cons.
On the same $80$ images of cohort~B with four runs each, it gains $4.9$ points over fixed best-of-$K$ at equal mean runs for DAPS, whose failed runs on these images are $12\%$ \cons, and $-0.1$ points for SITCOM, whose failed runs are $60\%$ \cons\claim{C-G2}.

\paragraph{Development results.}
Runs of an image are exchangeable in our data\claim{C-F4}, so random orders of an image's runs simulate sequential sampling.
On a panel of $n$ runs of which $L\ge M$ pass, the stopping time averages $M(n+1)/(L+1)$ over orders, and the first $M$ passing runs form a uniform $M$-subset of the $L$, so the protocol's value has the same exact subset form as the baselines of \cref{sec:selection}; it agrees with our simulated evaluation to within $3\times10^{-4}$\claim{C-T13}.
For FAST-DIPS on FFHQ with $64$ runs per image (\cref{fig:teaser}b), the protocol with $N_{\max}=64$ returns a \good image for $91.3\%$ of the $180$ images at $8.4$ runs per image, against $89.1\%$ for fixed best-of-$K$ at the same mean number of runs and $91.2\%$ for best-of-$16$\claim{C-G1}.
It beats fixed best-of-$K$ at all $12$ operating points we tried ($M=1,\dots,4$, $\tau_S$ scaled by $0.997$, $1$ and $1.003$), by $0.8$ to $4.7$ points, and stays close to the fixed-$K$ ground-truth orbit oracle at equal mean compute.
$\textsc{flag}_C$ is raised for $8$ images, $6$ of the $7$ that FAST-DIPS never solves; the seventh yields \cons runs\claim{C-G1}\claim{C-F5}.
These numbers are in-sample: we designed the protocol on these runs.

\paragraph{Frozen held-out test.}
We froze the protocol, its parameters and four hypotheses before any held-out per-image result existed\claim{C-G1}, on three panels: FAST-DIPS on ImageNet ($100$ images, $16$ runs, $\tau_S$ transferred from FFHQ), SITCOM ($16$ runs) and DAPS ($8$ runs) on the $180$ FFHQ images.
H1: at equal mean runs, the protocol loses at most one point against fixed best-of-$K$ on each panel.
H2: it gains at least two points on DAPS, and more than on SITCOM and on ImageNet, whose failed runs are mostly \cons.
H3: at least half of the $\textsc{flag}_C$ images are unsolved on DAPS and on ImageNet.
H4: on ImageNet, outputs without either flag are \good at least five points more often than all outputs without $\textsc{flag}_C$.
\Cref{tab:protocol} reports the outcome.

\begin{table}[t]
\centering\footnotesize
\setlength{\tabcolsep}{1.3pt}
\begin{tabular}{@{}lccccccc@{}}
\toprule
 & & \multicolumn{2}{c}{Protocol} & \multicolumn{2}{c}{Fixed best-of-$K$} & & \\
\cmidrule(lr){3-4}\cmidrule(lr){5-6}
Panel & $N_{\max}$ & \good & Runs & eq.\ runs & $K{=}N_{\max}$ & $\Oorb$ & $\textsc{flag}_C$ \\
\midrule
FAST, FFHQ$^\dagger$ & 64 & 91.3 & 8.4 & 89.1 & 93.9 & 96.1 & 8 (6) \\
FAST, ImageNet & 16 & \multicolumn{6}{c}{\pending{H1, H3, H4}} \\
SITCOM, FFHQ & 16 & \multicolumn{6}{c}{\pending{H1, H2}} \\
DAPS, FFHQ & 8 & \multicolumn{6}{c}{\pending{H1--H3}} \\
\bottomrule
\end{tabular}
\caption{\textbf{The protocol ($M=2$) against fixed best-of-$K$.} \good: \% of images with a \good output; Runs: mean runs per image. Fixed best-of-$K$ uses $\Eone$ (with an orientation test) or $\Eo$ (without), at the protocol's mean runs and at $K=N_{\max}$, the runs available. $\Oorb$: ground-truth orbit oracle at $N_{\max}$. $\textsc{flag}_C$: flagged images (unsolved among them). $^\dagger$Development data, without the orientation step\claim{C-G1}.}
\label{tab:protocol}
\end{table}
